\documentclass[11pt,a4paper]{article}
\usepackage[margin=2.2cm]{geometry}
\usepackage{fontspec}
\usepackage{polyglossia}
\setmainlanguage{spanish}
\usepackage{booktabs}
\usepackage{longtable}
\usepackage{array}
\usepackage{hyperref}
\hypersetup{colorlinks=true,linkcolor=blue,urlcolor=blue}
\setlength{\parindent}{0pt}
\setlength{\parskip}{6pt}
\newcommand{\code}[1]{\texttt{\detokenize{#1}}}

\title{Informe final\\\large GA8-220501096-AA1-EV02 — Módulos integrados}
\author{Diego Alberto Parra Garzón\\Iguaque QR}
\date{17 de septiembre de 2026}

\begin{document}
\maketitle

\section{Resumen}
Iguaque QR es una aplicación Android para crear, personalizar, guardar y compartir códigos QR mediante un flujo guiado de seis pasos. La copia de entrega contiene el código fuente, configuración Gradle, documentación técnica, avisos de terceros, licencia GPLv3 y el ejecutable debug de pruebas.

\section{Resultado técnico}
La compilación \code{assembleDebug} y el análisis \code{lintDebug} terminaron correctamente. Las pruebas unitarias existentes terminan con \code{NO-SOURCE}, porque el proyecto no contiene casos automatizados implementados.

El módulo residual de cámara y escaneo fue retirado. Se conservaron ZXing Core, Room, el historial de QR generados, el guardado y el uso compartido. No se modificó la aplicación publicada ni la firma release.

\section{Validación en dispositivo}
Se verificó la variante \code{com.diegodebian.iguaqueqr.test} en un Xiaomi M2101K6G con Android 13/API 33. La aplicación publicada permaneció instalada. El usuario confirmó visualmente que la barra inferior aparece con ``Crear QR'' e ``Historial'', que los botones del asistente siguen visibles, que se puede regresar al creador y que el QR generado aparece en el historial.

\section{Ejecutable}
\begin{tabular}{>{\bfseries}p{4cm}p{10cm}}
Archivo & \code{app/build/outputs/apk/debug/app-debug.apk}\\
Paquete & \code{com.diegodebian.iguaqueqr.test}\\
Versión & \code{1.0.1-test}\\
SHA-256 & \code{EFA6D69F4307A113C3AE4799240C0E682E8D17C26C09D1D0A0FAA18909F6841C}
\end{tabular}

\section{Módulos, entradas y salidas}
\begin{longtable}{p{3.2cm}p{4.4cm}p{5.4cm}}
\toprule
\textbf{Módulo} & \textbf{Entrada} & \textbf{Salida}\\
\midrule
Bienvenida & Apertura de la app & Acceso al creador\\
Contenido & Texto o URL & Contenido confirmado\\
Personalización & Forma, marco, colores y corrección & Configuración QR\\
Legibilidad & Colores y contraste & Estado de legibilidad\\
Renderizado & Configuración confirmada & Bitmap QR\\
Resultado & Bitmap generado & QR listo\\
Guardado/compartir & Acción del usuario & PNG local o selector Android\\
Historial & Registro generado & Lista local persistente\\
\bottomrule
\end{longtable}

\section{Persistencia y dependencias}
Room utiliza la base local \code{qr_art_database}, tabla \code{qr_history}, versión de esquema 1. La entidad y el DAO se conservaron sin cambios estructurales. \code{generateQR()} inserta registros con \code{type = "generated"}. ZXing Core 3.5.2 participa en la generación; las demás dependencias y licencias se detallan en \code{THIRD_PARTY_NOTICES.md}.

\section{Pruebas}
\begin{longtable}{p{2.1cm}p{9.5cm}p{2.2cm}}
\toprule
\textbf{Caso} & \textbf{Verificación} & \textbf{Estado}\\
\midrule
CP-01 & Apertura de la aplicación & Aprobado\\
CP-02 & Pantalla informativa inicial & Aprobado\\
CP-16 & Generación del QR & Aprobado\\
CP-21 & QR generado visible en el historial & Aprobado\\
NAV-01 & Barra inferior visible & Aprobado\\
NAV-02 & Solo Crear QR e Historial & Aprobado\\
NAV-03 & Regreso a Crear QR & Aprobado\\
NAV-04 & Botones del asistente visibles & Aprobado\\
\bottomrule
\end{longtable}
Los demás casos de la matriz requieren registro manual individual y no se marcan como aprobados por inferencia. El escaneo interno y la cámara están fuera del alcance del producto.

\section{Requisitos de la evidencia}
\begin{longtable}{p{8.5cm}p{4.8cm}}
\toprule
\textbf{Requisito} & \textbf{Estado}\\
\midrule
Código fuente & Disponible en la copia limpia\\
Archivo ejecutable & APK debug generado\\
Repositorio de control de versiones & Pendiente de inicialización/publicación autorizada\\
Documentación de módulos, entradas y salidas & Disponible\\
Pruebas y ambientes & Parcialmente documentados\\
Configuración de base de datos & Room documentado; esquema versión 1\\
Manual técnico & Disponible\\
URL & Pendiente de definir o publicar\\
Informe final PDF & Este documento\\
\bottomrule
\end{longtable}

\section{Observaciones y congelamiento}
No hay pruebas unitarias implementadas. Permanecen documentos UX históricos que mencionan escaneo y deben excluirse o actualizarse antes de una entrega documental definitiva. Los archivos sensibles y configuraciones locales no forman parte de la copia limpia.

El desarrollo funcional queda congelado. Cualquier cambio posterior requiere una autorización nueva y separada.

\end{document}
